\documentclass[10pt,a4paper]{amsart}
\usepackage{amsmath,amssymb,amsthm,graphicx}
% For commutative diagrams we use Xy-pic
\usepackage[all]{xy}

% We use 2cell for 2-commutative diagrams.
\xyoption{2cell}
% Theorem environments.
%
\theoremstyle{plain}
\newtheorem{theorem}[subsection]{Theorem}
\newtheorem{proposition}[subsection]{Proposition}
\newtheorem{lemma}[subsection]{Lemma}

\theoremstyle{definition}
\newtheorem{definition}[subsection]{Definition}
\newtheorem{example}[subsection]{Example}
\newtheorem{exercise}[subsection]{Exercise}
\newtheorem{situation}[subsection]{Situation}

\theoremstyle{remark}
\newtheorem{remark}[subsection]{Remark}
\newtheorem{remarks}[subsection]{Remarks}

\numberwithin{equation}{subsection}

% Macros
%
\def\lim{\mathop{\mathrm{lim}}\nolimits}
\def\colim{\mathop{\mathrm{colim}}\nolimits}
\def\Spec{\mathop{\mathrm{Spec}}}
\def\Hom{\mathop{\mathrm{Hom}}\nolimits}
\def\Ext{\mathop{\mathrm{Ext}}\nolimits}
\def\SheafHom{\mathop{\mathcal{H}\!\mathit{om}}\nolimits}
\def\SheafExt{\mathop{\mathcal{E}\!\mathit{xt}}\nolimits}
\def\Sch{\mathit{Sch}}
\def\Mor{\mathop{\mathrm{Mor}}\nolimits}
\def\Ob{\mathop{\mathrm{Ob}}\nolimits}
\def\Sh{\mathop{\mathit{Sh}}\nolimits}
\def\NL{\mathop{N\!L}\nolimits}
\def\CH{\mathop{\mathrm{CH}}\nolimits}
\def\proetale{{pro\text{-}\acute{e}tale}}
\def\etale{{\acute{e}tale}}
\def\QCoh{\mathit{QCoh}}
\def\Ker{\mathop{\mathrm{Ker}}}
\def\Im{\mathop{\mathrm{Im}}}
\def\Coker{\mathop{\mathrm{Coker}}}
\def\Coim{\mathop{\mathrm{Coim}}}

% Boxtimes
%
\DeclareMathSymbol{\boxtimes}{\mathbin}{AMSa}{"02}

%
% Macros for moduli stacks/spaces
%
\def\QCohstack{\mathcal{QC}\!\mathit{oh}}
\def\Cohstack{\mathcal{C}\!\mathit{oh}}
\def\Spacesstack{\mathcal{S}\!\mathit{paces}}
\def\Quotfunctor{\mathrm{Quot}}
\def\Hilbfunctor{\mathrm{Hilb}}
\def\Curvesstack{\mathcal{C}\!\mathit{urves}}
\def\Polarizedstack{\mathcal{P}\!\mathit{olarized}}
\def\Complexesstack{\mathcal{C}\!\mathit{omplexes}}
% \Pic is the operator that assigns to X its picard group, usage \Pic(X)
% \Picardstack_{X/B} denotes the Picard stack of X over B
% \Picardfunctor_{X/B} denotes the Picard functor of X over B
\def\Pic{\mathop{\mathrm{Pic}}\nolimits}
\def\Picardstack{\mathcal{P}\!\mathit{ic}}
\def\Picardfunctor{\mathrm{Pic}}
\def\Deformationcategory{\mathcal{D}\!\mathit{ef}}
\usepackage{cleveref}
\crefname{originallemma}{original lemma}{original lemmas}
\Crefname{originallemma}{Original lemma}{Original lemmas}
\makeatletter
\newlabel{lemma-chow-ring-projective-bundle}{{02TX}{source}{}{sourcepointer-1}{}}
\newlabel{lemma-chow-ring-projective-bundle@cref}{{[originallemma][0][]02TX}{[1][0][]source}{}{}{}}
\makeatother
\crefname{originallemma}{original lemma}{original lemmas}
\Crefname{originallemma}{Original lemma}{Original lemmas}
\makeatletter
\newlabel{lemma-flat-pullback-proper-pushforward}{{02RG}{source}{}{sourcepointer-2}{}}
\newlabel{lemma-flat-pullback-proper-pushforward@cref}{{[originallemma][0][]02RG}{[1][0][]source}{}{}{}}
\makeatother
\crefname{originallemma}{original lemma}{original lemmas}
\Crefname{originallemma}{Original lemma}{Original lemmas}
\makeatletter
\newlabel{lemma-pushforward-cap-c1}{{02SU}{source}{}{sourcepointer-3}{}}
\newlabel{lemma-pushforward-cap-c1@cref}{{[originallemma][0][]02SU}{[1][0][]source}{}{}{}}
\makeatother
\crefname{originallemma}{original lemma}{original lemmas}
\Crefname{originallemma}{Original lemma}{Original lemmas}
\makeatletter
\newlabel{situation-setup}{{02QL}{source}{}{sourcepointer-4}{}}
\newlabel{situation-setup@cref}{{[originallemma][0][]02QL}{[1][0][]source}{}{}{}}
\makeatother
\title{Characterisation of the classes c\_j of a finite locally free sheaf by a single relation on its projective bundle}
\author{The Stacks Project Authors}
\date{stacks-project:chow.tex:local-snapshot-20261012; GNU FDL 1.2 or later}
\begin{document}
\maketitle
\noindent\textit{Editorial note.} All delivered mathematical text is the author's own, copied from the pinned Stacks Project chapter file chow.tex; no statement or proof step was rewritten, reordered or removed.\par\medskip
\noindent\textit{Editorial note.} Preamble provenance: the pinned chapter file opens with \textbackslash{}input\{preamble\} and carries no preamble text of its own, so the project's single Stacks-wide preamble staged at sources/stacks-chow/preamble.tex is excerpted verbatim as two spans, lines 14-18 (Xy-pic setup) and lines 154-217 (the native theorem declarations and mathematical macros). Nothing was added, removed or redefined; the wrapper class and the amsmath/amssymb/amsthm/graphicx base packages are supplied by the assembler.\par\medskip
\noindent\textit{Editorial note.} Same-chapter pointers are located at their \textbackslash{}label lines inside the pinned chapter file; cross-chapter pointers are located in the project's own tag list staged in the snapshot (sources/stacks-chow/tags.txt), at the verbatim lines carrying their canonical tags.\par\medskip
\noindent\textit{Allowed standard prerequisites (editorial).}\begin{itemize}
\item The projective space bundle formula. Same chapter chow.tex, label lemma-chow-ring-projective-bundle, invoked at proof line 7303. Gives the uniqueness of alpha\_0, ..., alpha\_r with alpha\_0 = alpha satisfying the displayed relation, by identifying the Chow group of the projective bundle with the direct sum of the Chow groups of the base.
\item Flat pullback versus proper pushforward, and the projection formula. Same chapter chow.tex, labels lemma-flat-pullback-proper-pushforward (proof line 7348) and lemma-pushforward-cap-c1 (proof line 7349). Allows the relation proved on the disjoint union of the projective bundles over the components to be transported back to P, and then to the relation with pi\textasciicircum{}*(p\_* beta\_j) on X.
\end{itemize}\medskip
\ifdefined\hypertarget\hypertarget{sourcepointer-1}{}\fi
\noindent\textit{Source reference.} Original lemma 02TX: sources/stacks-chow/pinned-original-source.tex, source lines 6932--6932. Navigation pointer for \textbackslash{}ref\{lemma-chow-ring-projective-bundle\} used inside the counted span (line 7303); the argument that consumes it is retained in the span.\par\medskip
\ifdefined\hypertarget\hypertarget{sourcepointer-2}{}\fi
\noindent\textit{Source reference.} Original lemma 02RG: sources/stacks-chow/pinned-original-source.tex, source lines 2401--2401. Navigation pointer for \textbackslash{}ref\{lemma-flat-pullback-proper-pushforward\} used inside the counted span (line 7348); the argument that consumes it is retained in the span.\par\medskip
\ifdefined\hypertarget\hypertarget{sourcepointer-3}{}\fi
\noindent\textit{Source reference.} Original lemma 02SU: sources/stacks-chow/pinned-original-source.tex, source lines 4631--4631. Navigation pointer for \textbackslash{}ref\{lemma-pushforward-cap-c1\} used inside the counted span (line 7349); the argument that consumes it is retained in the span.\par\medskip
\ifdefined\hypertarget\hypertarget{sourcepointer-4}{}\fi
\noindent\textit{Source reference.} Original lemma 02QL: sources/stacks-chow/pinned-original-source.tex, source lines 1372--1372. Navigation pointer for \textbackslash{}ref\{situation-setup\} used inside the counted span (line 7281); the argument that consumes it is retained in the span.\par\medskip
\begin{lemma}
\label{lemma-determine-intersections}
Let $(S, \delta)$ be as in Situation \ref{situation-setup}.
Let $X$ be locally of finite type over $S$.
Let $\mathcal{E}$ be a finite locally free sheaf of rank $r$ on $X$.
Let $(\pi : P \to X, \mathcal{O}_P(1))$ be the projective bundle
associated to $\mathcal{E}$.
For $\alpha \in Z_k(X)$ the elements
$c_j(\mathcal{E}) \cap \alpha$ are the unique elements
$\alpha_j$ of $\CH_{k - j}(X)$
such that $\alpha_0 = \alpha$ and
$$
\sum\nolimits_{i = 0}^r
(-1)^i c_1(\mathcal{O}_P(1))^i \cap
\pi^*(\alpha_{r - i}) = 0
$$
holds in the Chow group of $P$.
\end{lemma}

\begin{proof}
The uniqueness of $\alpha_0, \ldots, \alpha_r$ such that
$\alpha_0 = \alpha$ and such that
the displayed equation holds follows from
the projective space bundle formula
Lemma \ref{lemma-chow-ring-projective-bundle}.
The identity holds by definition for $\alpha = [W]$ where $W$
is an integral closed subscheme of $X$.
For a general $k$-cycle $\alpha$ on $X$ write
$\alpha = \sum n_a[W_a]$ with $n_a \not = 0$, and
$i_a : W_a \to X$ pairwise distinct integral closed subschemes.
Then the family $\{W_a\}$ is locally finite on $X$.
Set $P_a = \pi^{-1}(W_a) = \mathbf{P}(\mathcal{E}|_{W_a})$.
Denote $i'_a : P_a \to P$ the corresponding closed immersions.
Consider the fibre product diagram
$$
\xymatrix{
P' \ar@{=}[r] \ar[d]_{\pi'} &
\coprod P_a \ar[d]_{\coprod \pi_a} \ar[r]_{\coprod i'_a} &
P \ar[d]^\pi \\
X' \ar@{=}[r] &
\coprod W_a \ar[r]^{\coprod i_a} &
X
}
$$
The morphism $p : X' \to X$ is proper. Moreover
$\pi' : P' \to X'$ together with the invertible sheaf
$\mathcal{O}_{P'}(1) = \coprod \mathcal{O}_{P_a}(1)$
which is also the pullback of $\mathcal{O}_P(1)$
is the projective bundle associated to
$\mathcal{E}' = p^*\mathcal{E}$. By definition
$$
c_j(\mathcal{E}) \cap [\alpha]
=
\sum i_{a, *}(c_j(\mathcal{E}|_{W_a}) \cap [W_a]).
$$
Write $\beta_{a, j} = c_j(\mathcal{E}|_{W_a}) \cap [W_a]$
which is an element of $\CH_{k - j}(W_a)$. We have
$$
\sum\nolimits_{i = 0}^r
(-1)^i c_1(\mathcal{O}_{P_a}(1))^i \cap \pi_a^*(\beta_{a, r - i}) = 0
$$
for each $a$ by definition. Thus clearly we have
$$
\sum\nolimits_{i = 0}^r
(-1)^i c_1(\mathcal{O}_{P'}(1))^i \cap (\pi')^*(\beta_{r - i}) = 0
$$
with $\beta_j = \sum n_a\beta_{a, j} \in \CH_{k - j}(X')$. Denote
$p' : P' \to P$ the morphism $\coprod i'_a$.
We have $\pi^*p_*\beta_j = p'_*(\pi')^*\beta_j$
by Lemma \ref{lemma-flat-pullback-proper-pushforward}.
By the projection formula of Lemma \ref{lemma-pushforward-cap-c1}
we conclude that
$$
\sum\nolimits_{i = 0}^r
(-1)^i c_1(\mathcal{O}_P(1))^i \cap \pi^*(p_*\beta_j) = 0
$$
Since $p_*\beta_j$ is a representative of $c_j(\mathcal{E}) \cap \alpha$
we win.
\end{proof}

\end{document}
